\label{sec:reproducibility}

To ensure transparency and enable independent verification of our research findings, we provide comprehensive resources for reproducing the spatial contagion simulation framework. This section details the availability of source code, dependencies, and step-by-step instructions for replicating our experiments.

\subsection{Source Code Availability}

The complete source code for the ChatDev-based contagion simulation platform is publicly available on GitHub:
\begin{itemize}
    \item \textbf{Repository}: \url{https://github.com/LaansDole/multi-agent-simulation}
    \item \textbf{Upstream Foundation}: \url{https://github.com/OpenBMB/ChatDev}
\end{itemize}

The repository contains all implementation artifacts, including the agent orchestration framework, cellular automata visualization module, and epidemiological model configurations. The codebase is released under the Apache-2.0 license, permitting both academic and commercial use with appropriate attribution.

\subsection{Dependencies and Environment Setup}

\subsubsection{System Requirements}

The simulation platform requires the following system specifications:
\begin{itemize}
    \item \textbf{Operating System}: macOS, Linux, Windows (via WSL)
    \item \textbf{Python}: Version 3.12 (not compatible with Python 3.13 or higher)
    \item \textbf{Node.js}: Version 18 or higher
    \item \textbf{Package Manager}: uv (recommended) or pip
    \item \textbf{Memory}: Minimum 8 GB RAM (16 GB recommended for complex scenarios)
\end{itemize}

\subsubsection{System Dependencies}

Before installing Python packages, ensure the following system libraries are installed:

\textbf{macOS (via Homebrew):}
\begin{verbatim}
brew install cairo pkg-config
\end{verbatim}

\textbf{Linux (Debian/Ubuntu):}
\begin{verbatim}
sudo apt-get update
sudo apt-get install -y libcairo2-dev pkg-config
\end{verbatim}

\subsubsection{Python Dependencies}

The project uses \texttt{uv} for dependency management. Core dependencies include:
\begin{itemize}
    \item \texttt{pyyaml}: YAML configuration parsing
    \item \texttt{openai}: LLM API integration
    \item \texttt{fastapi} and \texttt{uvicorn}: Web server framework
    \item \texttt{pydantic}: Data validation
    \item \texttt{networkx}: Graph-based workflow orchestration
    \item \texttt{matplotlib}, \texttt{seaborn}, \texttt{cartopy}: Visualization
    \item \texttt{pygame}: Interactive simulation rendering
\end{itemize}

The complete dependency list is specified in \texttt{pyproject.toml} at the repository root.

\subsection{Installation and Configuration}

\subsubsection{Step 1: Clone the Repository}
\begin{verbatim}
git clone https://github.com/LaansDole/multi-agent-simulation.git
cd multi-agent-simulation
\end{verbatim}

\subsubsection{Step 2: Install Backend Dependencies}
\begin{verbatim}
uv sync
\end{verbatim}

\subsubsection{Step 3: Install Frontend Dependencies}
\begin{verbatim}
cd frontend && npm install && cd ..
\end{verbatim}

\subsubsection{Step 4: Configure Environment Variables}
\begin{verbatim}
cp .env.example .env
\end{verbatim}

Edit \texttt{.env} to set your LLM provider credentials:
\begin{verbatim}
BASE_URL=https://api.openai.com/v1
API_KEY=your_api_key_here
\end{verbatim}

\subsection{Running the Simulation}

\subsubsection{Option A: Web Console (Recommended)}

Start both backend and frontend services:
\begin{verbatim}
make dev
\end{verbatim}

Access the web interface at \url{http://localhost:5173}. From the Launch tab, select a workflow configuration and enter your simulation scenario.

\subsubsection{Option B: Command Line Interface}

Execute a workflow directly from the terminal:
\begin{verbatim}
uv run python run.py \
    --path yaml_instance/simulation_hospital.yaml \
    --name my_simulation
\end{verbatim}

\subsubsection{Option C: Python SDK}

For programmatic execution and integration:
\begin{verbatim}
from runtime.sdk import run_workflow

result = run_workflow(
    yaml_file="yaml_instance/simulation_hospital.yaml",
    task_prompt="COVID-19 outbreak in a 200-bed hospital",
    variables={"API_KEY": "sk-xxxx"}
)

print(result.final_message.text_content())
\end{verbatim}

\subsection{Sample Configuration Files}

We provide pre-configured workflow templates demonstrating different contagion scenarios:

\begin{itemize}
    \item \textbf{Hospital Simulation}: \texttt{yaml\_instance/simulation\_hospital.yaml}
    \begin{itemize}
        \item Multi-agent hospital environment with nurses, doctors, and patients
        \item Demonstrates disease transmission in healthcare settings
        \item Includes memory-based patient tracking and environmental awareness
        \item Example prompt: ``March 2020, COVID-19 outbreak beginning. Hospital scrambling to set up isolation protocols.''
    \end{itemize}
    
    \item \textbf{Workflow Template}: \texttt{yaml\_template/design.yaml}
    \begin{itemize}
        \item Base template for creating custom simulation scenarios
        \item Configurable agent roles, interaction patterns, and environmental parameters
        \item Suitable for school, workplace, or community contagion modeling
    \end{itemize}
\end{itemize}

Additional demo configurations (e.g., \texttt{demo\_*.yaml}) showcase specific platform features and can be adapted for educational contagion scenarios.

\subsection{Validation and Testing}

To verify correct installation and validate workflow configurations:

\begin{verbatim}
# Validate all YAML workflow files
make validate-yamls

# Run the test suite
uv run pytest

# Check specific workflow syntax
uv run python -m check.check --path yaml_instance/simulation_hospital.yaml
\end{verbatim}

\subsection{Persistent Identifier}

Upon publication, a DOI will be assigned to the release version of this software through Zenodo, enabling permanent citation and version tracking. The current preprint is available on arXiv~\cite{qian2023chatdev}.

\subsection{Reproducibility Statement}

Following the principles outlined by Munaf{\`o} et al.~\cite{munafo2017manifesto}, we commit to:
\begin{enumerate}
    \item \textbf{Open Data}: All simulation configurations and example outputs are included in the repository.
    \item \textbf{Version Control}: Git history preserves all code changes with descriptive commit messages.
    \item \textbf{Documentation}: Comprehensive README, inline comments, and user guides are provided.
    \item \textbf{Dependency Pinning}: Exact package versions are specified in \texttt{pyproject.toml}.
    \item \textbf{Containerization}: Docker support is available for environment reproducibility.
\end{enumerate}

Researchers seeking to extend or replicate our work should refer to the Developer Guide in \texttt{docs/user\_guide/} for detailed architecture documentation and extension patterns.
